% !TeX encoding = UTF-8
\documentclass{../plantilla/hvtbook}
\HVTTitulo{Orgánicos, microbiología y fermentación}
\HVTSubtitulo{Aprende a ver los residuos orgánicos como recurso: separación, microorganismos, fermentación oscura y rutas biológicas para energía y bioinsumos.}
\HVTNivel{Nivel aplicado}\HVTSerie{Guía formativa}\HVTNumero{04}\HVTAccion{Empezar}
\HVTDescripcion{Guía formativa sobre residuos orgánicos, microbiología y fermentación.}
\HVTCanonical{https://hidrogenoverdeturquesa.com/curso-organicos-fermentacion}
\HVTRegreso{Volver al catálogo}{/fundacion\#cursos}
\HVTPortada{../../images/Fermentacion Oscura.png}{Figura 1. La fermentación oscura aprovecha materia orgánica bajo condiciones controladas.}
\begin{document}\HVTMakeTitle
\begin{HVTPage}{Capítulo 1}{Residuo o recurso}{Objetivo de la lección}
\HVTLead{Al terminar, el estudiante clasifica residuos orgánicos, reconoce el papel de los microorganismos y entiende la fermentación como una ruta de transformación biológica.}
\begin{HVTGrid}
\begin{HVTColumn}{Aprenderás}\begin{itemize}\item Qué residuos orgánicos se pueden estudiar. \item Qué necesitan los microorganismos para trabajar. \item Qué significa ambiente sin oxígeno. \item Cómo registrar cambios de olor, color, pH y gas.\end{itemize}\end{HVTColumn}
\begin{HVTColumn}{Materiales de estudio}\begin{itemize}\item Lista de residuos de cocina. \item Frascos para demostración cerrada no presurizada. \item Etiquetas y cuaderno. \item Guantes y normas de higiene.\end{itemize}\end{HVTColumn}
\end{HVTGrid}
\HVTImagen{../../images/Fermentacion Oscura.png}{Figura 1-1. Un sistema biológico debe observarse con orden, higiene y registro de variables.}
\end{HVTPage}
\begin{HVTPage}{Concepto}{Microorganismos como transformadores}{La idea central}
\HVTText{Los microorganismos transforman materia orgánica en nuevos compuestos. En fermentación oscura, algunas rutas pueden producir gases y ácidos orgánicos bajo condiciones controladas sin luz y con poco o ningún oxígeno.}
\begin{HVTFlow}\HVTFlowItem{1}{Sustrato}{Residuo orgánico seleccionado y preparado.}\HVTFlowItem{2}{Microbiología}{Comunidad de microorganismos transforma el sustrato.}\HVTFlowItem{3}{Productos}{Gases, líquidos, ácidos y material estabilizado.}\end{HVTFlow}
\begin{HVTExample}{Ejemplo 4-1}{Clasificación de residuos para una práctica}\begin{enumerate}\item Separar frutas y vegetales frescos. \item Evitar mezclas con aceites, químicos o contaminantes. \item Registrar masa inicial y fecha. \item Observar cambios sin abrir innecesariamente.\end{enumerate}\end{HVTExample}
\HVTText{Comentario: la calidad del sustrato define la estabilidad del proceso. Separar bien es una decisión técnica, no solo ambiental.}
\end{HVTPage}
\begin{HVTPage}{Práctica}{Diario de observación}{Registra antes de intervenir}
\begin{HVTSteps}\item Selecciona el residuo: define origen, tipo y fecha. \item Pesa o estima cantidad: registra masa o volumen inicial. \item Etiqueta el recipiente: nombre, fecha y responsable. \item Observa variables: color, olor, burbujas, pH si hay indicador disponible. \item Compara: anota cambios diarios y posibles causas.\end{HVTSteps}
\begin{HVTCode}{Formato de registro}día:
sustrato:
olor:
color:
gas observado:
pH o indicador:
comentario:\end{HVTCode}
\end{HVTPage}
\begin{HVTPage}{Actividad}{Evidencia de aprendizaje}{Mapa de aprovechamiento en casa}
\HVTText{Diseña un ecosistema doméstico donde los orgánicos no terminan en basura sino en compost, bioinsumos, estudio microbiológico o energía.}
\begin{HVTTask}{Entrega}\item Dibuja el flujo de orgánicos de una casa. \item Marca qué residuos van a compost y cuáles a estudio. \item Escribe tres reglas de higiene y seguridad. \item Explica cómo esto se conecta con una casa energética.\end{HVTTask}
\end{HVTPage}
\end{document}
